\begin{slikaDesno}{fig/R_ser.pdf}
\PID \label{rser}
Два отпорника чије су измерене отпорности са одговарајућим \textit{апсолутним} грешкама 
$R_1 = (670 \pm 30)\unit{\Omega}$ и 
$R_2 = (820 \pm 40)\unit{\Omega}$, повезани су редно као на слици. 
Одредити еквивалентну отпорност те редне везе $R_{\rm e}$, као и \textit{апсолутну} 
грешку ове отпорности $\Delta R_{\rm e}$ 
\end{slikaDesno}

Отпорност редне везе је 
\begin{equation}
    R_{\rm e} = R_1 + R_2 = 1490\unit{\Omega}. \label{eq:\ID:1}
\end{equation}

Израз за апсолутну грешку мотивисан је тоталним диференцијалном фукције више променљивих, односно, 
уколико је за $y = y(x_1, x_2, \ldots, x_N)$, онда је тотални диференцијал дат у облику 
$
\de y = \dfrac{\partial y}{\partial x_1} \, \de x_1 + \cdots + \dfrac{\partial y}{\partial x_N} \, \de x_N.
$
Приликом процене \textit{апсолутне грешке} рачуна се најгори случај, однсоно да се сви доприноси сабирају, чиме 
се узима апсолутна вредност сваког од чланова. Додатно, сматра се да су грешке довољно мале, да се 
израз за тотални диференцијал може искористити за процену грешке мерења, односно да се заменом 
$\de x_i \mapsto \Delta x_i$ добија процена апсолутне грешке целокупног израза $\de y \mapsto \Delta y$ као 
\begin{equation}
    \Delta y = \left| \dfrac{\partial y}{\partial x_1} \right| \Delta x_1 +
    \cdots + \left| \dfrac{\partial y}{\partial x_N} \right| \Delta x_N 
    = \sum_{i = 1}^N \left(\left| \dfrac{\partial y}{\partial x_i} \right| \Delta x_i  \right)
    \label{eq:def_aps_greska}
\end{equation}
Применом овог резултата на израз \eqref{eq:\ID:1} се онда добија 
\begin{equation}
    \Delta R_{\rm e} 
    = 
    \cancelto{1}{\left| \dfrac{\partial R_{\rm e}}{\partial R_1}  \right|}
    \Delta R_1
    +
    \cancelto{1}{\left| \dfrac{\partial R_{\rm e}}{\partial R_2}  \right|}
    \Delta R_2 
    = \Delta R_1 + \Delta R_2 = 70\unit{\Omega}
\end{equation}
Коначно, резултат за отпорност редне везе се може записати у облику\footnote{По правилу, ако није другачије 
наглашено, грешка се мајорира на једну значајну цифру, а резултат се заокругљује на истом броју 
сигурних цифара.} 
$R_{\rm e} = (1490 \pm 70)\unit{\Omega}$. 